{\songti\fontsize{9.5pt}{13.5pt}\selectfont
\noindent\textbf{摘要：}以溶胶-凝胶法和静电纺丝法分别制备 TiO$_2$ 前驱体，并在空气气氛中经 200—900 $^\circ$C 保温 4 h 热处理，采用 X 射线衍射（XRD）、拉曼光谱、扫描电子显微镜与能谱（SEM/EDS）以及紫外-可见漫反射光谱对产物进行系统表征。结果表明：200 $^\circ$C 样品近于非晶态，仅呈现宽化的衍射包；300—600 $^\circ$C 产物为锐钛矿相，其 (101) 晶面衍射峰半高宽由 1.19$^\circ$ 收窄至 0.29$^\circ$，Scherrer 晶粒尺寸由 6.8 nm 增至 28.1 nm；700 $^\circ$C 起出现金红石相（质量分数 17.3\%），800 $^\circ$C 升至 84.7\%，900 $^\circ$C 完全转变为金红石。拉曼光谱中锐钛矿的 144、399、516、639 cm$^{-1}$ 特征峰随温度升高逐渐减弱，金红石的 448 和 611 cm$^{-1}$ 峰逐渐增强，与 XRD 结果一致；700 $^\circ$C 两个平行样品分别呈纯锐钛矿与纯金红石，说明相变在微区上并不均匀。SEM 统计得到 300、500、700 $^\circ$C 颗粒平均尺寸分别为 31$\pm$11、39$\pm$12、85$\pm$28 nm，900 $^\circ$C 时晶粒粗化至亚微米级；静电纺丝产物为直径中位数约 0.48 $\mu$m 的连续纤维，并伴有少量珠状缺陷。EDS 仅检出 Ti、O、C 及来自样品台的 Al，未发现其他杂质。由 Tauc 关系外推得到 400—600 $^\circ$C 样品的表观禁带宽度为 3.12—3.13 eV，700 $^\circ$C 降至 3.05 eV，800 $^\circ$C 进一步降至 2.80 eV。上述结果说明热处理温度通过“相变—晶粒生长—形貌粗化”的耦合过程同时调控 TiO$_2$ 的结构与光学性质，而制备路线主要影响产物的宏观组装形态。\par}
